\chapter*{Lời Mở Đầu}
\addcontentsline{toc}{chapter}{Lời Mở Đầu}
\markboth{Lời Mở Đầu}{Lời Mở Đầu}

\begin{truyen}{Lời Mở Đầu}{Opening Words}
\chuhoa{B}{ộ sách này bắt đầu từ một buổi kể chuyện đơn giản trên lớp.}
\textit{(This series began from a simple storytelling session in class.)}

Thầy giáo kể, học sinh nghe, rồi cả lớp đặt câu hỏi. Và thầy nhận ra rằng những câu chuyện nhỏ lại chứa những bài học lớn — đủ lớn để một học sinh mang theo suốt đời.
\textit{(The teacher told stories, students listened, and the whole class asked questions. The teacher realized that small stories carry large lessons — large enough for a student to carry for a lifetime.)}

Quyển đầu tiên này là lời chào của cả bộ sách. Những điển tích ở đây nói về lòng biết ơn, sự can đảm, tình bạn nghĩa khí và ý chí vươn lên.
\textit{(This first volume is the greeting of the entire series. The stories here speak of gratitude, courage, loyal friendship, and the will to rise.)}

Mỗi câu chuyện kèm song ngữ Việt -- Anh để người đọc vừa học ý, vừa học tiếng.
\textit{(Each story comes with bilingual Vietnamese -- English so readers learn both meaning and language.)}
\end{truyen}

\begin{baihoc}
Những câu chuyện hay nhất không cần phải dài --- chúng chỉ cần chạm đúng vào điều người ta đang cần nghe.
\textit{(The best stories do not need to be long --- they only need to touch exactly what people need to hear.)}
\end{baihoc}

\begin{ghinhoanh}
\textbf{Short English to remember:}

Good stories carry timeless lessons.

Stories teach what lectures cannot.
\end{ghinhoanh}

\begin{tuhocnhanh}
1. Điều gì trong những câu chuyện này khiến em suy nghĩ nhất? \textit{Which story in this volume made you think the most?}

2. Bài học nào em muốn áp dụng ngay hôm nay? \textit{Which lesson do you want to apply today?}

3. Nếu kể lại một câu chuyện cho người khác, em chọn câu nào? \textit{If you were to retell one story to others, which would you choose?}
\end{tuhocnhanh}

\ngancach
